\documentclass[exercice]{thedocument}%
\usepackage{texmaths-tikz}%
\startdatakeys
\source{}
\cycle{}
\competencesDuSocle{}
\connaissancesRequises{}
\competencesTravaillees{}
\enddatakeys
\begin{document}%
    \begin{center}
        \begin{tikzpicture}
            \coordinate[label=below left:\<s1;](s1)at(0,0);
            \coordinate[label=below right:\<s2;](s2)at(4,0);
            \coordinate[label=above right:\<s3;](s3)at(4,2.5);
            \coordinate[label=above left:\<s4;](s4)at(0,2.5);
            \tkzMarkRightAngle(s1,s2,s3)
            \tkzMarkRightAngle(s2,s3,s4)
            \tkzMarkRightAngle(s3,s4,s1)
            \tkzMarkRightAngle(s4,s1,s2)
            \draw(s1)--(s2)--(s3)--(s4)--cycle;
        \end{tikzpicture}
    \end{center}
    \begin{enumerate}
        \item En répondant par des phrases, indiquer quelles droites sont
        perpendiculaires sur la figure ci-dessus.
        \item Traduire chacune des phrases de la question \textbf{1)} en
        utilisant les notations mathématiques.
    \end{enumerate}
\end{document}
